% ═══════════════════════════════════════════════════════════════
% SPANISCH ARBEITSBLATT: Verbos en -ar (Repaso) im Schulkontext
% ═══════════════════════════════════════════════════════════════
% Thema: -ar Verben wiederholen + Schulalltag
% Klasse 7 | Unidad 2 | Mi instituto | DS 04
% ═══════════════════════════════════════════════════════════════

\documentclass[10pt, a4paper]{article}

\usepackage[utf8]{inputenc}
\usepackage[T1]{fontenc}
\usepackage[ngerman]{babel}

% --- SCHRIFTEN ---
\usepackage{palatino}
\usepackage{helvet}
\newcommand{\sanstext}[1]{{\fontfamily{phv}\selectfont #1}}

\usepackage{microtype}
\usepackage[margin=1.4cm, top=1.6cm, bottom=1.3cm]{geometry}
\usepackage{enumitem}
\usepackage{tabularx}
\usepackage{booktabs}
\usepackage{array}
\usepackage{multicol}
\usepackage{tikz}
\usetikzlibrary{calc}

\usepackage{fancyhdr}
\usepackage{lastpage}
\usepackage{needspace}

\usepackage{xcolor}
\usepackage{amssymb}
\usepackage{tcolorbox}
\tcbuselibrary{skins}

% ═══════════════════════════════════════════════════════════════
% FARBEN
% ═══════════════════════════════════════════════════════════════

\definecolor{kopfzeile}{HTML}{1A1A1A}
\definecolor{akzent}{HTML}{C41E3A}          % Spanisch-Karminrot
\definecolor{hellgrau}{HTML}{F2F2F2}
\definecolor{mittelgrau}{HTML}{777777}
\definecolor{dunkelgrau}{HTML}{444444}

% ═══════════════════════════════════════════════════════════════
% VARIABLEN
% ═══════════════════════════════════════════════════════════════

\newcommand{\thema}{Verbos en -ar (Repaso)}
\newcommand{\klasse}{7}
\newcommand{\maxpunkte}{30}

% ═══════════════════════════════════════════════════════════════
% KOPF- UND FUSSZEILE
% ═══════════════════════════════════════════════════════════════

\pagestyle{fancy}
\fancyhf{}
\renewcommand{\headrulewidth}{0pt}
\renewcommand{\footrulewidth}{0pt}
\cfoot{\sanstext{\footnotesize\textcolor{mittelgrau}{\thepage\,/\,\pageref{LastPage}}}}

% ═══════════════════════════════════════════════════════════════
% BOXEN
% ═══════════════════════════════════════════════════════════════

% Grammatik-Box
\newtcolorbox{grammatikbox}{
    colback=hellgrau,
    colframe=hellgrau,
    boxrule=0pt,
    arc=0pt,
    left=8pt, right=8pt, top=8pt, bottom=6pt,
    borderline north={1.5pt}{0pt}{dunkelgrau}
}

% Text-Box (für Lesetext)
\newtcolorbox{textbox}{
    colback=white,
    colframe=mittelgrau,
    boxrule=0.5pt,
    arc=0pt,
    left=10pt, right=10pt, top=8pt, bottom=8pt
}

% Merk-Box
\newtcolorbox{merkbox}{
    colback=hellgrau,
    colframe=hellgrau,
    boxrule=0pt,
    arc=0pt,
    left=8pt, right=8pt, top=8pt, bottom=6pt,
    borderline north={1.5pt}{0pt}{akzent}
}

% ═══════════════════════════════════════════════════════════════
% BEFEHLE
% ═══════════════════════════════════════════════════════════════

% Spanisch: Palatino fett
\newcommand{\es}[1]{\textbf{#1}}

% Deutsch: klein, kursiv, grau
\newcommand{\de}[1]{{\small\textit{\textcolor{mittelgrau}{#1}}}}

% Lücke
\newcommand{\luecke}[1]{\rule{#1}{0.4pt}}

% Punktebox
\newcommand{\punktebox}[1]{\hfill\sanstext{\small[\_\_/#1]}}

% Aufgabentitel
\newcommand{\aufgabe}[2]{%
    \needspace{4\baselineskip}%
    \vspace{10pt}%
    \noindent\sanstext{\textbf{#1}\quad#2}%
    \vspace{5pt}%
    \nopagebreak[4]%
}

% Teilaufgaben
\newlist{teil}{enumerate}{1}
\setlist[teil]{label=\alph*), leftmargin=1.8em, itemsep=4pt, parsep=0pt, topsep=3pt}

\setlength{\parindent}{0pt}
\setlength{\parskip}{4pt}
\setlength{\columnsep}{24pt}

% ═══════════════════════════════════════════════════════════════
% DOKUMENT
% ═══════════════════════════════════════════════════════════════

\begin{document}

% ═══════════════════════════════════════════════════════════════
% HEADER
% ═══════════════════════════════════════════════════════════════

\noindent
\begin{tikzpicture}[remember picture, overlay]
    \fill[kopfzeile] (current page.north west) rectangle +(\paperwidth, -1cm);
    \node[anchor=west, text=white, font=\sanstext\large\bfseries]
        at ([xshift=1.4cm, yshift=-0.5cm]current page.north west) {\thema};
    \node[anchor=east, text=white, font=\sanstext\small]
        at ([xshift=-1.4cm, yshift=-0.5cm]current page.north east)
        {Spanisch\,·\,Klasse \klasse\,·\,Arbeitsblatt};
\end{tikzpicture}

\vspace{0.5cm}

% --- Name/Datum ---
\noindent
\sanstext{\small\textbf{Name:}} \luecke{5cm} \hfill
\sanstext{\small\textbf{Datum:}} \luecke{2.5cm}

\vspace{8pt}

% ═══════════════════════════════════════════════════════════════
% GRAMMATIK-ÜBERSICHT
% ═══════════════════════════════════════════════════════════════

\begin{grammatikbox}
\begin{multicols}{2}

% --- Konjugation ---
{\large\bfseries Verbos en -ar} \de{Konjugation}

\vspace{3pt}

{\small
\begin{tabular}{@{}ll@{\hspace{12pt}}ll@{}}
yo & habl\es{o} & nosotros & habl\es{amos} \\
tú & habl\es{as} & vosotros & habl\es{áis} \\
él/ella & habl\es{a} & ellos/ellas & habl\es{an} \\
\end{tabular}}

\columnbreak

% --- Verben-Liste ---
{\large\bfseries Verbos importantes}

\vspace{3pt}

{\footnotesize
hablar \de{sprechen} · estudiar \de{lernen} · escuchar \de{zuhören} \\
mirar \de{schauen} · buscar \de{suchen} · tomar \de{nehmen} \\
descansar \de{ausruhen} · charlar \de{plaudern} · esperar \de{warten}}

\end{multicols}
\end{grammatikbox}

\vspace{2pt}

% ═══════════════════════════════════════════════════════════════
% AUFGABEN
% ═══════════════════════════════════════════════════════════════

\begin{multicols}{2}

% --- AUFGABE 1 ---
\aufgabe{1}{Konjugiere die Verben.} \punktebox{6}

\vspace{4pt}

{\small
\renewcommand{\arraystretch}{1.3}
\begin{tabular}{@{}l|p{1.6cm}|p{1.6cm}|p{1.6cm}@{}}
 & \textbf{estudiar} & \textbf{escuchar} & \textbf{mirar} \\
\hline
yo & \luecke{1.4cm} & \luecke{1.4cm} & \luecke{1.4cm} \\
tú & \luecke{1.4cm} & \luecke{1.4cm} & \luecke{1.4cm} \\
él/ella & \luecke{1.4cm} & \luecke{1.4cm} & \luecke{1.4cm} \\
nosotros & \luecke{1.4cm} & \luecke{1.4cm} & \luecke{1.4cm} \\
ellos & \luecke{1.4cm} & \luecke{1.4cm} & \luecke{1.4cm} \\
\end{tabular}}

\vspace{8pt}

% --- AUFGABE 2 ---
\aufgabe{2}{Ergänze das passende Verb.} \punktebox{6}

{\small Benutze: \es{hablar, estudiar, escuchar, mirar, buscar, descansar}}

\vspace{4pt}

\begin{teil}
    \item En clase de español nosotros \luecke{2cm} español.

    \vspace{2pt}

    \item Los alumnos \luecke{2cm} la pizarra.

    \vspace{2pt}

    \item Yo \luecke{2cm} matemáticas todos los días.

    \vspace{2pt}

    \item María \luecke{2cm} su bolígrafo. ¿Dónde está?

    \vspace{2pt}

    \item Vosotros \luecke{2cm} al profesor.

    \vspace{2pt}

    \item En el recreo los alumnos \luecke{2cm}.
\end{teil}

\columnbreak

% --- AUFGABE 3: Lesetext ---
\aufgabe{3}{Lee el texto y contesta.} \punktebox{6}

\begin{textbox}
{\small
Hola, me llamo Pablo. \es{Soy} alumno en el Instituto García Lorca. \es{Tengo} catorce años.

Por la mañana llego al instituto a las ocho. En mi clase \es{hay} veinticinco alumnos. La profesora de español se llama señora Martínez.

En la primera clase \es{estudio} matemáticas. \es{Escucho} al profesor y tomo apuntes. Después \es{tengo} clase de español. \es{Hablamos} español y \es{miramos} la pizarra.

En el recreo \es{descanso} y \es{charlo} con mis amigos. \es{Buscamos} un banco en el patio.}
\end{textbox}

\vspace{4pt}

{\small Beantworte auf Spanisch:}

\begin{teil}
    \item ¿Cuántos alumnos hay en la clase de Pablo?

    \luecke{5cm}

    \vspace{2pt}

    \item ¿Qué hace Pablo en clase de español? (2 cosas)

    \luecke{5cm}

    \vspace{2pt}

    \item ¿Qué hacen Pablo y sus amigos en el recreo?

    \luecke{5cm}
\end{teil}

\end{multicols}

% ═══════════════════════════════════════════════════════════════
% AUFGABE 4: Volle Breite
% ═══════════════════════════════════════════════════════════════

\vspace{2pt}
\noindent\textcolor{hellgrau}{\rule{\textwidth}{1.5pt}}
\vspace{2pt}

\aufgabe{4}{Bilde Sätze im Schulkontext.} \punktebox{6}

{\small Verbinde die Elemente und schreibe vollständige Sätze. Benutze auch \es{hay/estar}!}

\vspace{4pt}

\begin{multicols}{3}
{\footnotesize
\textbf{Subjekt:}\\
yo, tú, María, nosotros, los alumnos, el profesor}

\columnbreak

{\footnotesize
\textbf{Verb (-ar):}\\
hablar, estudiar, escuchar, mirar, buscar, tomar}

\columnbreak

{\footnotesize
\textbf{Objekt/Ort:}\\
español, la pizarra, al profesor, apuntes, el libro, en la clase}
\end{multicols}

\vspace{4pt}

1. \luecke{14cm}

\vspace{6pt}

2. \luecke{14cm}

\vspace{6pt}

3. \luecke{14cm}

% ═══════════════════════════════════════════════════════════════
% AUFGABE 5: Produktion
% ═══════════════════════════════════════════════════════════════

\vspace{6pt}

\aufgabe{5}{Mi día en el instituto – Schreibe einen kurzen Text.} \punktebox{6}

{\small Schreibe 3-4 Sätze über deinen Schulalltag. Benutze Verben auf \es{-ar} und auch \es{ser, tener, estar, hay}.}

\vspace{4pt}

{\footnotesize\textit{Hilfe:} Soy alumno/alumna... Tengo ... años. En mi clase hay... Estudio... Hablo... Escucho...}

\vspace{6pt}

\noindent\fbox{\parbox{\dimexpr\textwidth-2\fboxsep-2\fboxrule}{%
\vspace{2.5cm}
}}

% ═══════════════════════════════════════════════════════════════
% MERKKASTEN
% ═══════════════════════════════════════════════════════════════

\vspace{6pt}

\begin{merkbox}
\sanstext{\textbf{Merke: Verben im Schulkontext}}

\vspace{3pt}

\begin{multicols}{2}

{\small
\begin{tabular}{@{}ll@{}}
\es{hablar} español & Spanisch sprechen \\
\es{estudiar} matemáticas & Mathe lernen \\
\es{escuchar} al profesor & dem Lehrer zuhören \\
\end{tabular}}

\columnbreak

{\small
\begin{tabular}{@{}ll@{}}
\es{mirar} la pizarra & die Tafel anschauen \\
\es{buscar} el libro & das Buch suchen \\
\es{tomar} apuntes & Notizen machen \\
\end{tabular}}

\end{multicols}
\end{merkbox}

% ═══════════════════════════════════════════════════════════════
% FOOTER
% ═══════════════════════════════════════════════════════════════

\vfill

\noindent\rule{\textwidth}{0.5pt}

\vspace{4pt}

\hfill\sanstext{\textbf{Punkte:}} \luecke{1.2cm} / \maxpunkte

\end{document}
